\documentclass[11pt]{article}

\usepackage[margin=1in]{geometry}
\usepackage{amsmath,amssymb,amsthm}
\usepackage{graphicx}
\usepackage{booktabs}
\usepackage{microtype}
\usepackage[dvipsnames]{xcolor}
\usepackage[colorlinks=true,linkcolor=NavyBlue,citecolor=NavyBlue,urlcolor=NavyBlue]{hyperref}
\usepackage{caption}
\usepackage{enumitem}

\captionsetup{font=small,labelfont=bf}
\graphicspath{{figures/}}

\newtheorem{theorem}{Theorem}
\newtheorem{proposition}{Proposition}
\newtheorem{corollary}{Corollary}
\theoremstyle{definition}
\newtheorem{definition}{Definition}
\theoremstyle{remark}
\newtheorem{remark}{Remark}

% Generated by paper/make_numbers.py -- do not edit by hand.
% Every measurement quoted in the paper is defined here from
% results/data/*.json, so the prose cannot drift from the data.

\newcommand{\EfiveAbsReffMax}{4.6}
\newcommand{\EfiveAbsReffMin}{2.3}
\newcommand{\EfiveDarkAUC}{0.561}
\newcommand{\EfiveDarkEff}{98.4}
\newcommand{\EfiveDomEff}{17.4}
\newcommand{\EfiveFisherAUC}{0.775}
\newcommand{\EfiveFisherEff}{9.2}
\newcommand{\EfiveKLgbar}{0.301}
\newcommand{\EfiveKLprobe}{0.019}
\newcommand{\EfiveMeanReff}{0.31}
\newcommand{\EfiveModels}{gpt2, distilgpt2, pythia-70m}
\newcommand{\EfiveN}{400}
\newcommand{\EfiveNfamilies}{2}
\newcommand{\EfiveNmodels}{3}
\newcommand{\EfiveSecond}{distilgpt2}
\newcommand{\EfiveSecondDarkEff}{98.8}
\newcommand{\EfiveSecondDomEff}{13.7}
\newcommand{\EfiveSecondFisherEff}{7.3}
\newcommand{\EfiveThirdDarkEff}{98.6}
\newcommand{\EfiveThirdDomEff}{14.1}
\newcommand{\EfiveThirdFisherEff}{8.2}
\newcommand{\EfiveThirdName}{pythia-70m}
\newcommand{\EfiveWidths}{512, 768}
\newcommand{\EfourAucNinety}{14.7}
\newcommand{\EfourAucNinetyList}{32/6/6}
\newcommand{\EfourAucNinetyNine}{15.3}
\newcommand{\EfourCertNinety}{1.3}
\newcommand{\EfourCertNinetyList}{2/1/1}
\newcommand{\EfourCertNinetyNine}{1.3}
\newcommand{\EfourLayers}{3, 6, 9}
\newcommand{\EfourM}{64}
\newcommand{\EfourNlayers}{3}
\newcommand{\EfourOracleNinety}{1.0}
\newcommand{\EfourOracleNinetyList}{1/1/1}
\newcommand{\EfourOracleNinetyNine}{1.0}
\newcommand{\EfourRandNinety}{29.3}
\newcommand{\EfourRandNinetyList}{46/22/20}
\newcommand{\EfourRandNinetyNine}{49.0}
\newcommand{\EfourSpAUC}{0.313}
\newcommand{\EfourSpSbar}{0.761}
\newcommand{\EfourSpeedupAUC}{11}
\newcommand{\EfourSpeedupRandom}{22}
\newcommand{\EoneAxisPopEffect}{$0$}
\newcommand{\EoneCorrP}{0.061}
\newcommand{\EoneCorrS}{0.124}
\newcommand{\EoneDarkAUC}{0.502}
\newcommand{\EoneDarkGap}{$5\times10^{-14}$}
\newcommand{\EoneEmpEffect}{$2.0\times10^{-2}$}
\newcommand{\EonePopEffect}{$2\times10^{-15}$}
\newcommand{\EoneProbeAUC}{0.998}
\newcommand{\EoneProbeAUCclosed}{0.998}
\newcommand{\EsixLateMatched}{0.062}
\newcommand{\EsixLateUnmatched}{0.039}
\newcommand{\EsixMatched}{0.031}
\newcommand{\EsixN}{300}
\newcommand{\EsixNbeh}{5}
\newcommand{\EsixNlab}{4}
\newcommand{\EsixNlayers}{12}
\newcommand{\EsixRatio}{0.97}
\newcommand{\EsixUnmatched}{0.032}
\newcommand{\EsixWins}{6}
\newcommand{\EthreeD}{768}
\newcommand{\EthreeDarkAUC}{0.676}
\newcommand{\EthreeDialLevels}{0.50, 0.70, 0.85, 0.95}
\newcommand{\EthreeDialRhos}{0.066, 0.051, 0.057, 0.047}
\newcommand{\EthreeEffDark}{99.7}
\newcommand{\EthreeEffProbe}{6.0}
\newcommand{\EthreeEffRandom}{2.6}
\newcommand{\EthreeFPRidentZero}{50}
\newcommand{\EthreeFPRneutZero}{0}
\newcommand{\EthreeGbarAUC}{0.698}
\newcommand{\EthreeHostileGap}{15.5}
\newcommand{\EthreeIdentAUClast}{0.520}
\newcommand{\EthreeInWarranty}{10}
\newcommand{\EthreeMaxAUC}{1.000}
\newcommand{\EthreeMeanReff}{0.18}
\newcommand{\EthreeMeanRho}{0.056}
\newcommand{\EthreeMinRho}{0.006}
\newcommand{\EthreeModel}{gpt2}
\newcommand{\EthreeNlayers}{12}
\newcommand{\EthreePairGap}{1.044}
\newcommand{\EthreePairN}{96}
\newcommand{\EthreePairSign}{71.9}
\newcommand{\EthreeSpAUC}{-0.063}
\newcommand{\EthreeSpSbar}{0.917}
\newcommand{\EtwoAbsReff}{1.6}
\newcommand{\EtwoAcc}{1.000}
\newcommand{\EtwoD}{64}
\newcommand{\EtwoDarkHi}{98.8}
\newcommand{\EtwoDarkLo}{98.7}
\newcommand{\EtwoDeltaInSD}{0.2}
\newcommand{\EtwoEffHi}{12.0}
\newcommand{\EtwoEffList}{12.0, 12.1, 12.0, 12.0}
\newcommand{\EtwoEffLo}{12.0}
\newcommand{\EtwoLayers}{5}
\newcommand{\EtwoLevels}{0.50, 0.70, 0.90, 0.98}
\newcommand{\EtwoMarkerList}{0.504, 0.703, 0.899, 0.979}
\newcommand{\EtwoNlevels}{4}
\newcommand{\EtwoNruns}{12}
\newcommand{\EtwoNseeds}{3}
\newcommand{\EtwoRhoDelta}{-0.004}
\newcommand{\EtwoRhoHi}{0.147}
\newcommand{\EtwoRhoList}{0.151, 0.150, 0.148, 0.147}
\newcommand{\EtwoRhoLo}{0.151}
\newcommand{\EtwoSeedSD}{0.017}
\newcommand{\PfiveMargin}{0.081}
\newcommand{\PfiveNrand}{400{,}000}
\newcommand{\PsevenAltRtwo}{0.721}
\newcommand{\PsevenAltSlope}{0.904}
\newcommand{\PsevenRtwo}{0.959}
\newcommand{\PsevenSlope}{1.006}
\newcommand{\PthreeIso}{0.975}
\newcommand{\PthreeMaxRelErr}{$1.3\times10^{-2}$}


\newcommand{\Sbar}{\bar{S}}
\newcommand{\gbar}{\bar{g}}
\newcommand{\Wg}{W_{g}}
\newcommand{\reff}{r_{\mathrm{eff}}}

\title{\textbf{Seeing Is Not Steering}\\[2mm]
\large Effect-Gramian Certificates for Causally Usable Latent Directions}

\author{Blessings Mambwe\\
\small \texttt{bleymambwe@gmail.com}}

\date{\today}

\begin{document}
\maketitle

\begin{abstract}
\noindent
Interpretability routinely fits a linear probe to a model's hidden states and
then reuses that same direction as a steering vector. This conflates two
different roles: the probe is an \emph{observation operator}, the steering
vector an \emph{actuator}. We show the two are governed by disjoint objects.
The probe direction is determined by the representation geometry
$(\Sigma_\ell,\delta_\ell)$; the steering-effective directions are determined by
the \emph{effect Gramian} $\Wg(\ell)=\mathbb{E}[gg^\top]$ over adjoint covectors
$g=\nabla_{x_\ell}\Phi_\ell$. Neither enters the other's definition, so their
alignment $\rho$ is a free parameter. We prove it can be set arbitrarily:
there is a system in which the optimal probe attains AUC $\ge 1-\epsilon$ and
steers \emph{exactly} nothing, while a chance-level direction steers maximally.
We introduce the \emph{steerability certificate}
$\Sbar(\ell,w)=|\gbar_\ell^\top w|/\|w\|$, computable from one reverse-mode pass
and \emph{zero} interventions, together with a closed form for probe-orthogonal
directions, an exact expression for the fraction of control an uninformed
direction retains ($\reff/d$), and an explicit validity radius. Empirically, on
\EthreeModel{} a toxicity probe classifying held-out prompts at AUC
\EthreeMaxAUC{} recovers only \EthreeEffProbe\% of the causal control available
at equal write norm, while a direction constrained to be orthogonal to it
recovers \EthreeEffDark\%. Only \EthreeMeanReff\% of the residual stream carries
first-order behavioural effect. In a controlled comparison, a probe trained on
\emph{exactly} the concept a behaviour measures is no better aligned with that
behaviour ($\rho=\EsixMatched$) than a probe for an unrelated concept
($\rho=\EsixUnmatched$; ratio \EsixRatio). The certificate rank-predicts
measured steering (Spearman \EthreeSpSbar) where probe accuracy does not
(\EthreeSpAUC), and using it to order interventions reaches 90\% of the best
attainable effect in \EfourCertNinety{} measurements out of \EfourM{} against
\EfourRandNinety{} for random selection. We report three results against
interest, including a pre-registered prediction that failed.
\end{abstract}

\section{Introduction}

A recurring pattern in mechanistic interpretability runs as follows. Collect
hidden states from a frozen language model, fit a linear probe until it detects
some concept---toxicity, sentiment, refusal, truthfulness---at high accuracy,
and conclude that the model ``represents'' that concept. Then take the probe
direction, add a multiple of it back into the residual stream, and expect the
model's behaviour to move accordingly.

The probe is being asked to do two jobs. As a \emph{sensor} it must separate
classes; as an \emph{actuator} it must move a behaviour. These are different
requirements, and nothing in the fitting procedure ties them together.

That this gap exists is increasingly acknowledged. What is missing is a theory
that \emph{predicts} it. Existing work observes the discrepancy after running
the interventions, which is precisely when the information is least useful. We
give the structural reason for the gap, a quantitative certificate that
anticipates it from cheap adjoint information, and measurements of its size.

\paragraph{Contributions.}
\begin{enumerate}[leftmargin=1.5em,itemsep=1pt]
  \item \textbf{A formalisation} of the frozen Transformer as an exact discrete
  controlled system, with probes as observation operators and steering as
  actuation, yielding the effect Gramian $\Wg$ and the certificate $\Sbar$
  (\S\ref{sec:setup}--\S\ref{sec:two}).
  \item \textbf{An exact decoupling theorem} (Thm.~\ref{thm:decouple}): probe
  accuracy and steering efficacy are independently settable. Corollary: no
  function of probe accuracy alone consistently estimates steering efficacy.
  \item \textbf{Structural results}: the fraction of control an uninformed
  direction retains is exactly $\reff/d$ (Prop.~\ref{prop:inert}); a closed form
  for the strongest probe-orthogonal direction (Prop.~\ref{prop:dark}); an
  explicit first-order validity radius (Prop.~\ref{prop:radius}); and a
  finite-sample leakage law $\Theta(N^{-1/2})$, \emph{independent of width}
  (Prop.~\ref{prop:leak}).
  \item \textbf{Measurements} on constructed systems, on transformers trained
  from scratch, and on three pretrained models with real toxicity data, using
  difference-of-means and logistic probes---the vectors practitioners actually
  deploy---as first-class baselines.
  \item \textbf{Active intervention design}: ordering candidate directions by
  the certificate reaches 90\% of the best attainable effect in
  \EfourCertNinety{} measurements out of \EfourM{}, against \EfourRandNinety{}
  for random selection --- a $\sim\EfourSpeedupRandom\times$ budget reduction
  (\S\ref{sec:results}).
  \item \textbf{Negative results}, including a pre-registered prediction that
  failed and a metric we discarded as vacuous (\S\ref{sec:negative}).
\end{enumerate}

\section{Setup: forward system, sensor, actuator}
\label{sec:setup}

Let a frozen decoder-only Transformer have residual stream $x_\ell\in\mathbb{R}^d$,
$\ell=0,\dots,L$, evolving as
\begin{equation}
x_{\ell+1} = x_\ell + F_\ell(x_\ell;c),
\label{eq:fwd}
\end{equation}
with $c$ a fixed context. Equation~\eqref{eq:fwd} is exact: no continuous-time,
autonomy, invertibility or energy-descent assumption is made or required.

Fix a scalar behaviour $\varphi:\mathbb{R}^d\to\mathbb{R}$ (in experiments, a
logit difference between a target and a contrast token set), and define the
layer-$\ell$ tail map
$\Phi_\ell := \varphi \circ T_{L-1}\circ\cdots\circ T_\ell$ where
$T_\ell(x)=x+F_\ell(x;c)$.

A linear probe at layer $\ell$ is a covector $w$ with readout $w^\top x_\ell$:
an observation operator $H_\ell=w^\top$. An additive intervention is
$x_\ell \mapsto x_\ell + \alpha w$, with measured causal effect
\begin{equation}
\Delta_i(\alpha,w) := \Phi_\ell\!\left(x_\ell^{(i)}+\alpha w\right)-\Phi_\ell\!\left(x_\ell^{(i)}\right).
\label{eq:effect}
\end{equation}

\begin{quote}
\textbf{Assumption A1.} $\Phi_\ell$ is $C^2$ on a neighbourhood of the segment
$[x_\ell^{(i)},x_\ell^{(i)}+\alpha w]$, with $\|\nabla^2\Phi_\ell\|_{\mathrm{op}}\le M_\ell$ there.
\end{quote}

A1 holds for GELU/SiLU blocks; $M_\ell$ is \emph{estimated} in
\S\ref{sec:results}, not assumed.

\section{Two geometries and a certificate}
\label{sec:two}

\begin{definition}[Adjoint covector and effect Gramian]
For example $i$, $g_\ell^{(i)} := \nabla_x\Phi_\ell(x_\ell^{(i)}) = J_\ell^{(i)\top}\nabla\varphi(x_L^{(i)})$,
obtained by one reverse-mode pass. Write
$\Wg(\ell) := \tfrac1N\sum_i g_\ell^{(i)}g_\ell^{(i)\top}$ and
$\gbar_\ell := \tfrac1N\sum_i g_\ell^{(i)}$.
\end{definition}

$\Wg$ is the finite-horizon, behaviour-restricted analogue of a
controllability-to-output Gramian for \eqref{eq:fwd}: it measures how much
\emph{behaviour} a unit-norm write at layer $\ell$ can produce. Its kernel,
$\mathcal{N}(\ell)=\ker\Wg(\ell)$, is inert to first order for every sampled example.

\begin{remark}[Non-normal composition as a candidate mechanism]
\label{rem:nonnormal}
By the chain rule,
\[
J_\ell^{(i)} = DT_{L-1}(x_{L-1}^{(i)})\cdots DT_\ell(x_\ell^{(i)}),
\]
a product of the per-block Jacobians $DT_k=I+DF_k(\cdot\,;c)$: this is exact, not
a limiting approximation. Each factor is generically non-normal---nothing forces
$DF_k$ to be symmetric---and products of non-normal matrices can have
singular-value spectra far more concentrated than the spectral radius of any
individual factor would suggest, the classical phenomenon of transient
amplification despite well-behaved eigenvalues~\cite{trefethen1993}. Direct
measurements on production-scale LLMs report exactly this signature in the
residual-stream Jacobian---a learned non-normality gradient through depth and a
cumulative low-rank perturbation bottleneck~\cite{nonnormalresidual2026}---and in
the attention propagator specifically, via pseudospectral and Kreiss-style
diagnostics~\cite{transientreserves2026}. This is a plausible structural
mechanism for the small $\reff(\ell)$ measured directly from $\Wg(\ell)$ in
\S\ref{sec:results}: a highly non-normal $J_\ell^{(i)}$ would concentrate
$g_\ell^{(i)}$, and hence $\Wg(\ell)$, onto a small subspace independently of
whether the underlying computation is itself low-dimensional. We do not test
this here---no Schur, pseudospectral, or Kreiss diagnostic is computed on these
models---and add it to \S\ref{sec:next}.
\end{remark}

On the read side, given binary labels $y_i$, let $\delta_\ell=\mu_1-\mu_0$ and
$\Sigma_\ell$ be the pooled within-class covariance. The Fisher-optimal probe is
$w_{\mathrm{probe}} \propto (\Sigma_\ell+\gamma I)^{-1}\delta_\ell$.

\begin{remark}[The source of the phenomenon]
$\Wg$ does not depend on $\Sigma_\ell$; $w_{\mathrm{probe}}$ does not depend on
$\Wg$. They are computed from disjoint information, and nothing in training
couples them. Their alignment
$\rho(\ell) := |\cos\angle(w_{\mathrm{probe}},\gbar_\ell)|$ is therefore a free
parameter of the system.
\end{remark}

\begin{definition}[Steerability certificate]
For $w\neq 0$,
\begin{equation}
\Sbar(\ell,w) := \frac{|\gbar_\ell^\top w|}{\|w\|},
\qquad
S_2(\ell,w) := \frac{\sqrt{w^\top \Wg(\ell)\, w}}{\|w\|}.
\end{equation}
\end{definition}

By Cauchy--Schwarz $\Sbar\le S_2\le\sqrt{\lambda_1(\Wg)}$, with $\Sbar$
maximised at $w=\gbar_\ell/\|\gbar_\ell\|$. Both are computable from $N$
reverse-mode passes over data the probe already required, and \emph{no
intervention is performed}.

\begin{theorem}[Three-way decomposition]
\label{thm:decomp}
Let $P_C$ project onto the top-$r$ eigenspace of $\Wg(\ell)$ with eigenvalues
$\lambda_1\ge\cdots\ge\lambda_d\ge0$, and $P_O$ onto the discriminative subspace
$\mathcal{O}(\ell)=\mathrm{span}\{(\Sigma_\ell+\gamma I)^{-1}\delta_\ell\}$. Then
for every $w$ and $i$, under A1,
\begin{equation}
\Delta_i(\alpha,w) = \alpha\left[g^{(i)\top}P_C w + g^{(i)\top}(I-P_C)w\right] + R_i,
\quad |R_i| \le \tfrac{M_\ell}{2}\alpha^2\|w\|^2,
\end{equation}
and $\mathrm{RMS}_i\,|g^{(i)\top}(I-P_C)w| \le \sqrt{\lambda_{r+1}}\|w\|$.
\end{theorem}

Consequently $\mathbb{R}^d$ partitions into \emph{causal} ($\mathcal{C}_r\cap\mathcal{O}$),
\emph{epiphenomenal} ($\mathcal{O}\setminus\mathcal{C}_r$: visible, inert),
\emph{dark} ($\mathcal{C}_r\setminus\mathcal{O}$: invisible, effective), and inert.
Membership is decided by two computable numbers, $\Sbar(\ell,w)$ and the probe
AUC of $w$. Standard practice assumes everything lies in the first cell.

\begin{remark}[Relation to the Kalman canonical decomposition]
\label{rem:kalman}
The four cells echo the controllable/observable canonical structure theorem
for linear time-invariant systems~\cite{kalman1963}, which partitions state
space into controllable-and-observable, controllable-only, observable-only,
and neither, using an exact controllability Gramian and an exact
observability Gramian built from the system matrices. The correspondence
here is structural, not literal: Kalman's theorem concerns the full state of
an autonomous linear system, while $\mathcal C_r(\ell)$ is a rank-$r$
truncation of a \emph{behaviour}-restricted, finite-horizon, first-order
Gramian, and $\mathcal O(\ell)$ is a \emph{label}-dependent discriminative
subspace with no counterpart in $(A,B,C)$. What transfers is the organising
idea---that ``can be moved'' and ``can be seen'' are governed by different
Gramians and need not coincide---not Kalman's exactness or invariance
results. Theorem~\ref{thm:decouple} is the certificate that this failure to
coincide can be total in the present, nonlinear, sampled setting.
\end{remark}

\section{Probe accuracy carries no causal information}

\begin{theorem}[Exact decoupling]
\label{thm:decouple}
For every $d\ge2$ and $\epsilon\in(0,\tfrac12)$ there is a system satisfying A1
---indeed with $\Phi_\ell$ linear, so $M_\ell=0$ and the first-order expansion is
exact---and a labelled representation such that
\begin{enumerate}[label=(\alph*),leftmargin=2em,itemsep=0pt]
  \item the optimal linear probe attains $\mathrm{AUC}\ge1-\epsilon$;
  \item $\Delta_i(\alpha,w_{\mathrm{probe}})=0$ \emph{exactly}, for all $\alpha$ and all $i$;
  \item there is a unit $v$ with $\mathrm{AUC}(v)=\tfrac12$ exactly and
  $|\gbar^\top v|=\|\gbar\|$, the maximum attainable.
\end{enumerate}
\end{theorem}

\begin{proof}[Proof sketch]
Take $\Phi_\ell(x)=e_1^\top x$, so $g^{(i)}=\gbar=e_1$; draw
$x\mid y\sim\mathcal{N}(y\Delta\mu\, e_2,\sigma^2 I_d)$, giving
$w_{\mathrm{probe}}\propto e_2$. Then (b) is
$\alpha e_1^\top e_2 = 0$; (a) follows from
$\mathrm{AUC}=\Phi_{\mathcal N}(\Delta\mu/(\sigma\sqrt2))$; and $v=e_1$ has a
class-independent score, giving (c). Full proof in Appendix~\ref{app:proofs}.
\end{proof}

\begin{corollary}
\label{cor:noinfo}
No function of probe accuracy alone is a consistent estimator of steering
efficacy: the pair $(\mathrm{AUC},\Sbar)$ can be set to any admissible values
independently.
\end{corollary}

Theorem~\ref{thm:decouple} also refutes the converse assumption---that a
probe-invisible direction is causally irrelevant---by construction.

\begin{proposition}[Random-direction inertness]
\label{prop:inert}
Let $\reff(\ell) := \operatorname{tr}\Wg(\ell)/\lambda_1(\Wg(\ell))$. For $w$
uniform on $S^{d-1}$,
\begin{equation}
\mathbb{E}\!\left[S_2(\ell,w)^2\right] = \frac{\operatorname{tr}\Wg(\ell)}{d},
\qquad
\frac{\mathbb{E}[S_2^2]}{\max_w S_2^2} = \frac{\reff(\ell)}{d}.
\end{equation}
\end{proposition}

So $\reff/d$ is exactly the share of steering power retained by a direction
chosen without write-side information. Measuring $\reff/d\ll1$ makes ``we found
a direction and it moved behaviour'' a claim about luck at that scale.

\begin{remark}[Relation to resolving power]
\label{rem:resolution}
$\mathcal N(\ell)=\ker\Wg(\ell)$ plays the role of the null space in
classical linear resolution theory~\cite{backusgilbert1968}: directions a
declared observation cannot see are assigned zero resolving power regardless
of the size of the underlying quantity. $\reff(\ell)$ is the corresponding
effective-rank count for $\Wg(\ell)$---the number of directions the
behaviour functional resolves at layer $\ell$---so $\reff/d$ is a
resolving-power fraction, not a claim about how many neurons ``matter.''
Backus and Gilbert's theory is posed for a linear forward operator mapping a
continuum model to finitely many exact data functionals; $\Wg(\ell)$ instead
comes from a first-order expansion of one behaviour functional restricted to
a finite sample. The transfer is in the concept, not the guarantees: a small
$\reff/d$ is evidence of low resolving power at this Gramian, not evidence
that the unresolved directions are inert to every future experiment.
Remark~\ref{rem:nonnormal} sketches one candidate mechanism for why the
resolving power measured here is so concentrated.
\end{remark}

\begin{proposition}[Probe suboptimality]
$\Sbar(\ell,w_{\mathrm{probe}}) = \|\gbar_\ell\|\cdot\rho(\ell)$, with equality
to the optimum iff $(\Sigma_\ell+\gamma I)^{-1}\delta_\ell \parallel \gbar_\ell$.
\end{proposition}

\begin{proposition}[Dark directions, closed form]
\label{prop:dark}
$v^\star := (I-P_O)\gbar_\ell / \|(I-P_O)\gbar_\ell\|$ maximises $\Sbar(\ell,v)$
over $v\perp\mathcal{O}(\ell)$, attaining $\|(I-P_O)\gbar_\ell\|$. Under
shared-covariance Gaussian class conditionals, any $v\perp\mathcal{O}(\ell)$ in
the whitened metric has probe AUC exactly $\tfrac12$.
\end{proposition}

This upgrades ``dark directions'' from an existence claim to a construction: one
projection of one adjoint mean.

\begin{proposition}[Validity radius]
\label{prop:radius}
Under A1, the certificate determines the sign and leading behaviour of
$\Delta_i(\alpha,w)$ whenever
$|\alpha| < \alpha^\star(\ell,w) := 2|\gbar_\ell^\top w|/(M_\ell\|w\|^2)$.
\end{proposition}

$M_\ell$ is estimated by second differences along the steering ray, so
$\alpha^\star$ is reported, not assumed.

\begin{proposition}[Finite-sample causal leakage]
\label{prop:leak}
Let the population probe be exactly inert, $\gbar^\top w_{\mathrm{pop}}=0$, and
let $\hat{w}$ be the Fisher probe from $N$ i.i.d.\ samples. Then
$\mathbb{E}[\hat{\rho}] = \Theta(N^{-1/2})$, with a constant that does
\emph{not} grow with $d$.
\end{proposition}

The dimension-cancellation is the practically important part: the artefact is
not mitigated by a wider residual stream. A probe fitted on $N\sim10^3$ carries
spurious alignment of order $3\times10^{-2}$ before any real effect exists.

\section{Experimental setup}
\label{sec:setup-exp}

All experiments run on \textbf{4 CPU cores, no GPU}. This constrains scale and
is stated wherever it matters.

\paragraph{E1 (exact verification).} Linear $\Phi$ with Gaussian
class-conditional representations, where the first-order theory is exact.
Evaluated both axis-aligned and conjugated by a random rotation, to exclude
coordinate artefacts.

\paragraph{E3 (\EthreeModel, toxicity).} Matched minimal pairs: hostile versus
friendly framing is the causal factor; identity-marked subjects are a confound
correlated with the label at $0.85$, as in real corpora. Only the frame differs
within a pair. $\varphi$ is a logit difference between mild insult tokens and
matched neutral tokens at the final position. Probes are fitted on a train split;
every AUC and every steering measurement is on held-out prompts. Write magnitude
is $\alpha=0.10\cdot\overline{\|x_\ell\|}$ at each layer.

\paragraph{E5 (real data, multiple models).} \EfiveN{} RealToxicityPrompts
prompts, balanced between high ($>0.6$) and low ($<0.05$) toxicity.
Difference-of-means (as deployed by activation-steering methods) and
logistic-regression probes are included as first-class baselines.

\paragraph{E2 (trained from scratch).} \EtwoNruns{} transformers
(\EtwoLayers{} layers, $d=\EtwoD$, LayerNorm + attention + GELU) trained on a
synthetic task with a causal content variable and a spurious surface marker,
across \EtwoNlevels{} confound levels $\times$ \EtwoNseeds{} seeds. This tests
whether the gap is produced by training itself, and whether training under a
confound widens it.

\paragraph{E4 (active design).} \EfourM{} candidate directions at layers
\EfourLayers{} --- random draws, probe/adjoint mixtures, effect eigenvectors and
sparse coordinate writes. Each is measured once to establish ground truth, then
four orderings are compared at every budget.

\paragraph{E6 (structural independence).} One model and \EsixN{} prompts held
fixed. \EsixNlab{} label definitions vary the read side; \EsixNbeh{} behaviour
functionals vary the write side. Some pairs are \emph{matched} by construction:
the probe is trained on exactly the concept the behaviour measures.

\paragraph{Instrument validation.} \texttt{auc}, \texttt{fisher\_probe},
\texttt{spearman} and \texttt{pearson} are pinned to scikit-learn and SciPy
references to $\sim10^{-16}$ (including heavy-tie and degenerate cases) in
\texttt{tests/test\_instruments.py}.

\section{Results}
\label{sec:results}

\subsection{The decoupling is exact, and the plane is fully occupied}

At the population probe, steering has effect \EonePopEffect{} (floating-point
zero) while classification is at AUC \EoneProbeAUC{}, matching the closed form
\EoneProbeAUCclosed{}. The constructed probe-orthogonal direction attains AUC
\EoneDarkAUC{} and the maximal effect, with gap \EoneDarkGap{}. The
axis-aligned construction gives \EoneAxisPopEffect{} exactly. Sweeping a family
of systems (Fig.~\ref{fig:decouple}) leaves probe AUC and the certificate
essentially uncorrelated: Pearson $r=\EoneCorrP$, Spearman $\EoneCorrS$.

Proposition~\ref{prop:inert} matches Monte-Carlo to within
\PthreeMaxRelErr{} relative error across construction ranks, with an
isotropic control at $\PthreeIso$ against an exact $1.0$. The closed-form dark
direction (Prop.~\ref{prop:dark}) beat the best of \PfiveNrand{} random
directions in the same subspace by a margin of $\PfiveMargin$.

\begin{figure}[t]
\centering
\begin{minipage}{0.48\textwidth}
\centering
\includegraphics[width=\linewidth]{fig1_decoupling.pdf}
\caption{Each point is one system; the class-mean direction is rotated away
from the effect direction by $\theta$ and separability is swept. Probe accuracy
and the certificate move independently (Cor.~\ref{cor:noinfo}).}
\label{fig:decouple}
\end{minipage}\hfill
\begin{minipage}{0.48\textwidth}
\centering
\includegraphics[width=\linewidth]{fig2_leakage.pdf}
\caption{Finite-sample leakage. True alignment is exactly zero; the estimate
decays as $N^{-1/2}$ with no dependence on width $d$
(Prop.~\ref{prop:leak}).}
\label{fig:leak}
\end{minipage}
\end{figure}

\subsection{Finite-sample probes manufacture causal effect}

Regressing $\log\hat\rho$ on $\log N^{-1/2}$ gives slope $\PsevenSlope$
(predicted $1.0$) with $R^2=\PsevenRtwo$. The dimension-dependent alternative
$\sqrt{d/N}$ fits worse (slope $\PsevenAltSlope$, $R^2=\PsevenAltRtwo$). The
dimension-free law wins (Fig.~\ref{fig:leak}).

\subsection{A perfect probe, nearly orthogonal to the mechanism}

On \EthreeModel, the behaviour functional is validated first: across
\EthreePairN{} exactly matched pairs it moves by $\EthreePairGap$ logits in the
correct direction on \EthreePairSign\% of pairs.

From layer~1 onward the probe separates held-out prompts at AUC
\EthreeMaxAUC{}. Its alignment with the mechanism averages $\rho=\EthreeMeanRho$
(minimum \EthreeMinRho), and only \EthreeMeanReff\% of the \EthreeD-dimensional
residual stream is behaviourally live (Fig.~\ref{fig:layerwise}).

Measured at equal write norm (Fig.~\ref{fig:effects}), the probe direction
recovers \EthreeEffProbe\% of the effect available to the certificate-optimal
direction; a random direction recovers \EthreeEffRandom\%; and a direction
constrained to be \emph{orthogonal to the probe} recovers \EthreeEffDark\%.

\begin{figure}[t]
\centering
\includegraphics[width=\textwidth]{fig3_gpt2_layerwise.pdf}
\caption{\textbf{Left:} probe AUC on held-out prompts against alignment $\rho$
with the adjoint direction. \textbf{Right:} the live fraction $\reff/d$.}
\label{fig:layerwise}
\end{figure}

\begin{figure}[t]
\centering
\includegraphics[width=\textwidth]{fig4_gpt2_effects.pdf}
\caption{Measured causal effect at equal write norm,
$\alpha=0.10\cdot\overline{\|x_\ell\|}$. The probe-orthogonal direction tracks
the adjoint almost exactly; the probe itself does not.}
\label{fig:effects}
\end{figure}

\paragraph{This is not ``gradients steer better''.}
That $\gbar$ moves $\varphi$ is true by construction and carries no information.
Three things do. First, the probe direction is what practitioners deploy, and
the size of its shortfall had not been measured. Second, the certificate
rank-predicts measured effect across a pool of directions including random
draws, probe/adjoint mixtures and effect eigenvectors---Spearman $\EthreeSpSbar$,
against $\EthreeSpAUC$ for probe accuracy---so it is a genuine predictor over
the space, not a restatement of one special case (Fig.~\ref{fig:calib}). Third,
the dissociation runs both ways: the best steering direction is itself a weak
detector (AUC $\EthreeGbarAUC$ against the probe's $\EthreeMaxAUC$).

\paragraph{Validity.} With $M_\ell$ estimated by second differences, the write
magnitude used lies inside the certificate's validity radius at
\EthreeInWarranty{} of \EthreeNlayers{} layers. Where it does not (the first two
layers) the first-order prediction is out of warranty and we say so rather than
quoting it.

\subsection{The two geometries are structurally independent}

E6 is the sharpest available test (Fig.~\ref{fig:indep}). A probe trained on
\emph{exactly} the concept a behaviour functional measures aligns with that
behaviour's adjoint at $\rho=\EsixMatched$. A probe for a completely unrelated
concept aligns at $\rho=\EsixUnmatched$. The ratio is \EsixRatio, and matched
beats unmatched in \EsixWins{} of \EsixNlayers{} layers---chance.

We state one nuance rather than averaging it away: over the final third of the
network matched alignment does exceed unmatched ($\EsixLateMatched$ versus
$\EsixLateUnmatched$). Both remain near zero---$\rho=0.06$ is roughly $86^\circ$
from the effect direction---so neither regime confers usable control. The claim
is not that matched probes carry literally no signal, but that the signal is too
small to act on and is swamped by which layer one happens to be at.

\begin{figure}[t]
\centering
\begin{minipage}{0.48\textwidth}
\centering
\includegraphics[width=\linewidth]{fig5_independence.pdf}
\caption{Alignment when the probe is trained on \emph{exactly} the concept the
behaviour measures (matched) versus an unrelated concept (unmatched). If the
read and write geometries were the same object, matched would dominate
everywhere.}
\label{fig:indep}
\end{minipage}\hfill
\begin{minipage}{0.48\textwidth}
\centering
\includegraphics[width=\linewidth]{fig7_calibration.pdf}
\caption{Predicted $|\alpha\Sbar|$, computed with no interventions, against
measured $|\Delta\varphi|$ for every (layer, direction) pair. The dashed line
is $y=x$; landing near it is stronger than a rank correlation.}
\label{fig:calib}
\end{minipage}
\end{figure}

\begin{figure}[t]
\centering
\includegraphics[width=0.52\textwidth]{fig6_realdata.pdf}
\caption{RealToxicityPrompts, two models. Every read-side construction
classifies well and steers poorly; a direction built orthogonal to all of them
does the reverse.}
\label{fig:real}
\end{figure}

\subsection{Real data and deployed steering vectors}

On RealToxicityPrompts (Fig.~\ref{fig:real}, Table~\ref{tab:real}), probe
accuracy sits at $\EfiveFisherAUC$---\emph{not} saturated---so the gap is not an
artefact of a maxed-out probe. The Fisher probe recovers \EfiveFisherEff\% of
available control; difference-of-means, the vector activation-steering methods
actually ship, recovers \EfiveDomEff\%; a direction orthogonal to all read-side
constructions recovers \EfiveDarkEff\% at AUC $\EfiveDarkAUC$. The pattern
replicates on \EfiveSecond{} (\EfiveSecondFisherEff\%, \EfiveSecondDomEff\%,
\EfiveSecondDarkEff\%) and, across a different architecture family, on
\EfiveThirdName{} (\EfiveThirdFisherEff\%, \EfiveThirdDomEff\%,
\EfiveThirdDarkEff\%).

\paragraph{A weak signal on how this scales.}
The three models span widths $\{\EfiveWidths\}$ and \EfiveNfamilies{}
architecture families. Across them the effect rank in \emph{absolute} terms is
small and roughly constant---$\reff$ between \EfiveAbsReffMin{} and
\EfiveAbsReffMax---rather than growing with $d$. If $\reff$ stays $O(1)$ as
width increases, then $\reff/d$ \emph{shrinks} with scale and the gap widens
rather than closing. We state this as suggestive and not established: it rests
on two distinct widths and three checkpoints, which is far too little to fit a
trend. It does, however, point the opposite way to the failure mode we
identify as the largest threat in \S\ref{sec:limits}.

\begin{table}[t]
\centering
\small
\caption{RealToxicityPrompts, averaged over all layers. Efficiency is measured
effect relative to the certificate-optimal direction at equal write norm; KL is
the mean divergence of the next-token distribution under the same write. Rows
are generated from \texttt{results/data/}, not transcribed.}
\label{tab:real}
% Generated by paper/make_numbers.py -- do not edit by hand.
\begin{tabular}{llrrrr}
\toprule
Model & Direction & Mean AUC & Mean $\rho$ & Efficiency & KL (nats) \\
\midrule
gpt2 & Fisher probe & 0.775 & 0.069 & 9.2\% & 0.019 \\
gpt2 & logistic probe & 0.821 & 0.100 & 12.9\% & --- \\
gpt2 & difference-of-means & 0.773 & 0.163 & 17.4\% & 0.022 \\
gpt2 & orthogonal to all three & 0.561 & 0.985 & 98.4\% & --- \\
gpt2 & adjoint $\gbar$ & 0.665 & 1.000 & 100\% & 0.301 \\
\midrule
distilgpt2 & Fisher probe & 0.798 & 0.053 & 7.3\% & 0.013 \\
distilgpt2 & logistic probe & 0.843 & 0.084 & 10.7\% & --- \\
distilgpt2 & difference-of-means & 0.758 & 0.145 & 13.7\% & 0.040 \\
distilgpt2 & orthogonal to all three & 0.553 & 0.988 & 98.8\% & --- \\
distilgpt2 & adjoint $\gbar$ & 0.652 & 1.000 & 100\% & 0.056 \\
\midrule
pythia-70m & Fisher probe & 0.750 & 0.067 & 8.2\% & 0.011 \\
pythia-70m & logistic probe & 0.810 & 0.103 & 13.6\% & --- \\
pythia-70m & difference-of-means & 0.771 & 0.114 & 14.1\% & 0.018 \\
pythia-70m & orthogonal to all three & 0.579 & 0.989 & 98.6\% & --- \\
pythia-70m & adjoint $\gbar$ & 0.678 & 1.000 & 100\% & 0.025 \\
\bottomrule
\end{tabular}

\end{table}

\subsection{The certificate buys a cheaper search}

If the certificate predicts effect, it should be usable to \emph{choose} which
interventions to run. E4 tests this directly: given \EfourM{} candidate
directions and a budget of $B$ intervention measurements, which ordering finds
the strongest direction fastest? Candidates include random draws,
probe/adjoint mixtures, effect eigenvectors and sparse coordinate writes;
ground truth is obtained by measuring all \EfourM{} once, offline.

Ordering by the certificate reaches 90\% of the best attainable effect after
\EfourCertNinety{} measurements on average (\EfourCertNinetyList{} across
layers \EfourLayers), against \EfourAucNinety{} for ordering by probe accuracy
(\EfourAucNinetyList) and \EfourRandNinety{} for random order
(\EfourRandNinetyList); the oracle needs \EfourOracleNinety{}. That is a
\textbf{$\sim\EfourSpeedupRandom\times$} reduction in intervention budget
against random selection and $\sim\EfourSpeedupAUC\times$ against probe
accuracy (Fig.~\ref{fig:budget}). Rank correlation with the true effect is
\EfourSpSbar{} for the certificate and \EfourSpAUC{} for $|\mathrm{AUC}-\tfrac12|$.

The certificate costs nothing extra here: it reuses adjoints already computed
for the probe-training data.

\begin{remark}[Relation to c-optimal experimental design]
\label{rem:elfving}
Ranking candidates by $\Sbar(\ell,w)=|\gbar_\ell^\top w|/\|w\|$ is a discrete
instance of \emph{c-optimal design}: choosing where to sample so as to best
estimate one linear functional $c^\top\theta$ of an unknown parameter, rather
than the whole parameter vector, is classically solved by Elfving's
theorem~\cite{elfving1952}, which characterises the optimal design point by
alignment with $c$ on the convex hull of the candidate set. Here $\gbar_\ell$
plays the role of $c$ and the \EfourM{} candidates play the role of the design
points, but two things do not transfer: the candidate pool is finite rather
than a continuum, and the task is a budgeted search for the single best point
rather than a continuous allocation, so Elfving's optimality and efficiency
guarantees do not apply as stated. What transfers is the reason a
\emph{single} functional is so much cheaper to design for than the whole
Gramian: c-optimality needs only alignment with one vector, where sensor
placement for the full state needs the Gramian's spectrum. That is also why
the certificate is cheap here---one backward pass, no eigendecomposition---while
a general-purpose observability audit is not.
\end{remark}

\begin{figure}[t]
\centering
\includegraphics[width=\textwidth]{fig8_budget.pdf}
\caption{Fraction of the best attainable effect discovered after $B$
intervention measurements, out of \EfourM{} candidates. The dotted line marks
90\%. Certificate ordering is close to the oracle; probe-accuracy ordering is
much closer to random.}
\label{fig:budget}
\end{figure}

\paragraph{Collateral damage, and a caveat about how to read it.}
Table~\ref{tab:real} also reports the KL divergence of the next-token
distribution under each write. At \emph{equal write norm} the adjoint direction
disturbs the distribution more than the probe does (0.301 versus 0.019 nats on
gpt2), which is unsurprising: it moves the behaviour more. Read naively, that
looks like a point against it.

It is not, because the comparison should be at \emph{equal effect}, not equal
norm. Behaviour change is first-order in the write magnitude while KL is
second-order (it is the Fisher--Rao metric to leading order), so scaling the
probe to match the adjoint's effect requires
$\alpha_{\mathrm{probe}}/\alpha_{\gbar}\approx 100/\EfiveFisherEff \approx 11$
and costs roughly $11^2$ times its measured KL --- about 2.3 nats, against
0.301 for the adjoint. On that basis the certificate-optimal direction is the
cheaper actuator by close to an order of magnitude. We flag this as an
extrapolation from the quadratic scaling rather than a measurement: verifying
it requires a matched-effect sweep, which we have not run.

\section{Results against interest}
\label{sec:negative}

\paragraph{A pre-registered prediction failed, in two independent designs.}
We predicted the gap would widen with the spurious correlation between a
surface confound and the label. It does not.

In the \emph{pretrained} setting (E3), sweeping
$P(\text{identity matches label})$ over $\{\EthreeDialLevels\}$ moves mean
$\rho$ over $\{\EthreeDialRhos\}$---no trend, and the gap is already near-total
at zero confounding. One might object that a pretrained model's representations
were fixed long before our confound existed, so the dial never had a chance to
act.

E2 removes that objection by controlling training itself. We train
\EtwoNruns{} transformers from scratch (\EtwoNseeds{} seeds $\times$
\EtwoNlevels{} confound levels, \EtwoLayers{} layers, $d=\EtwoD$) on a task with
a causal content variable and a spurious surface marker, to test accuracy
\EtwoAcc. The manipulation demonstrably works: accuracy against the marker
tracks the dial at $\{\EtwoMarkerList\}$. Yet across confound levels
$\{\EtwoLevels\}$, mean $\rho$ moves over $\{\EtwoRhoList\}$---a change of
$\EtwoRhoDelta$ end to end, which is $\EtwoDeltaInSD$ within-level seed standard
deviations (\EtwoSeedSD) and therefore indistinguishable from zero. Probe
steering efficiency is flat at $\{\EtwoEffList\}$\%, and the probe-orthogonal
direction holds at \EtwoDarkLo--\EtwoDarkHi\% throughout
(Fig.~\ref{fig:dial}).

Two experimental designs, one inheriting representations and one creating them,
both refute the confound hypothesis. This makes the result \emph{stronger} and
more general: the mechanism is the extreme low-rankness of $\Wg$
(Prop.~\ref{prop:inert}), not label contamination. Cleaning the labels will not
fix it.

\begin{figure}[t]
\centering
\includegraphics[width=\textwidth]{fig9_confound_dial.pdf}
\caption{The confound dial in both designs. \textbf{Left:} alignment $\rho$
against confound strength, for transformers trained from scratch (E2, error bars
over \EtwoNseeds{} seeds) and for pretrained gpt2 (E3). \textbf{Right:} probe
steering efficiency in E2. Neither shows the widening the pre-registered
prediction expected.}
\label{fig:dial}
\end{figure}

\paragraph{A proposition is only approximate in practice.}
Prop.~\ref{prop:dark} predicts probe-orthogonal directions classify at exactly
$0.500$. This is exact in the Gaussian construction (E1: \EoneDarkAUC) but not
in \EthreeModel, where they classify at $\EthreeDarkAUC$: real representations
are not shared-covariance Gaussian. The steering half of the proposition
survives intact (\EthreeEffDark\% of the optimum); the detectability half is
weaker than stated.

\paragraph{A metric was discarded as vacuous.} Our first collateral-damage
measure was the change in the prompt's own next-token perplexity under the
steer. It returned identically zero at every layer---structurally, not by a
small margin: the write lands on the final token, and causal masking forbids a
final-position state from influencing earlier predictions. We replaced it with
the KL divergence of the next-token distribution. We record this because a
metric that \emph{cannot} be non-zero is worse than no metric: it reads as
evidence of safety.

\paragraph{An identity-bias result that is smaller than it first appeared.}
On the confounded split the toxicity probe appears to detect identity mentions
at AUC up to $0.90$; on balanced $2\times2$ cells that figure falls to
$\EthreeIdentAUClast$ at the final layer. The disparity that persists is inside
the hostile class: hostile sentences with identity-marked subjects are flagged
\EthreeHostileGap{} percentage points more often than otherwise identical
sentences with neutral subjects. At layer $0$, benign identity-marked sentences
are flagged at \EthreeFPRidentZero\% against \EthreeFPRneutZero\%.

\section{Related work}

Circuit non-identifiability is now well established: multiple sufficient
subgraphs realise the same behaviour, and top-down and bottom-up analyses stand
in a many-to-many relation~\cite{eeao2025,uniq2026,manycircuits2026,certified2026}.
We take that as premise rather than contribution.

Control-theoretic framings of neural interpretability have recently
appeared---Gramian and Hankel-based mode ranking, and controllability analyses of
state-space language models~\cite{blackbox2025,ssmctrl2025}. Our contribution is
not the import of Gramians but their use as a \emph{predictive certificate for
probe-derived directions}, together with the closed form for probe-orthogonal
directions. The read/write split we formalise is a finite-horizon,
behaviour-restricted specialisation of three older results from outside
machine learning: Kalman's controllable/observable canonical decomposition
\cite{kalman1963}, which first separated ``can be moved'' from ``can be
seen'' as distinct subspaces of one linear system (Remark~\ref{rem:kalman});
Backus--Gilbert resolution theory~\cite{backusgilbert1968}, which quantifies
which directions a declared linear observation can and cannot resolve
(Remark~\ref{rem:resolution}); and Elfving's theory of c-optimal design
\cite{elfving1952} for a single linear functional, which the certificate
ordering in E4 specialises to a finite, budgeted candidate pool
(Remark~\ref{rem:elfving}). None of the three has previously been used, to
our knowledge, to certify or select individual steering directions from
adjoints already computed for probe training.

Independently, direct measurements of trained-Transformer Jacobians report a
learned non-normality gradient through depth and a resulting low-rank
perturbation bottleneck~\cite{nonnormalresidual2026,transientreserves2026},
the signature classical non-normal dynamical systems produce under
composition~\cite{trefethen1993}. Remark~\ref{rem:nonnormal} connects this to
the smallness of $\reff/d$ we measure directly from $\Wg$.

The gap between linear decodability and causal control is observed
qualitatively in recent steering work~\cite{langaxis2026}. Activation steering
and its difference-of-means constructions~\cite{actadd,caa}, feedback
controllers~\cite{feedback2025}, local-linearity optimal control~\cite{locallin2026}
and causal head selection for detoxification~\cite{causaldetox2026} all deploy
read-side directions as actuators. We supply the quantity that says in advance
how much of their effect is real. Probing methodology, activation and path
patching, ACDC and attribution patching, sparse autoencoders and causal
abstraction supply the surrounding toolkit; our object is orthogonal to circuit
discovery and concerns the validity of single directions.

\section{Limitations}
\label{sec:limits}

\textbf{Scale.} Everything ran on 4 CPU cores without a GPU; the largest model
is \EthreeModel{} (124M). The theory is scale-free, the measurements are not. If
$\reff/d$ grows with width, Prop.~\ref{prop:inert} loses its force; testing this
is the first item in \S\ref{sec:next}.

\textbf{First order.} The certificate is first-order.
Prop.~\ref{prop:radius} gives its validity radius, measured per layer, and we
mark where the prediction is out of warranty. Large-$\alpha$ steering studies
operate partly in that regime.

\textbf{One behaviour functional.} $\varphi$ is a logit difference over a small
token set. The minimal-pair check establishes that it tracks the intended
factor, not that it is the only reasonable readout.

\textbf{Position and pathway.} Interventions write to the final token's residual
stream at a single layer. Multi-position, multi-layer and attention-internal
writes are untested.

\textbf{What ``dark'' does not mean.} A dark direction is invisible to
\emph{this} probe family---linear, this concept, this data. It is not invisible
in principle; a nonlinear or differently-supervised detector may find it. The
safety claim concerns the monitors that are actually deployed.

\textbf{Toy-scale training.} The from-scratch models in E2 are small
(\EtwoLayers{} layers, $d=\EtwoD$). They establish that gradient descent
produces the gap unaided and that a training-time confound does not widen it,
but they are not evidence about what large-scale pretraining does.

\section{Safety implication}

If a monitor is built from a toxicity probe, it is close to blind on the
directions that most strongly drive toxic output: on \EthreeModel{} the
probe-orthogonal direction achieves \EthreeEffDark\% of the maximal effect while
the probe itself achieves \EthreeEffProbe\%. Probe-based monitoring is therefore
not, by itself, evidence of control. Given Theorem~\ref{thm:decouple}, this is a
consequence rather than a concern.

\section{Next steps}
\label{sec:next}

Ordered by what would most change the conclusions:
\begin{enumerate}[leftmargin=1.5em,itemsep=1pt]
  \item Replicate $\rho$ and $\reff/d$ at 1B--7B parameters. The single largest
  threat to external validity.
  \item Test Remark~\ref{rem:nonnormal} directly: compute Schur non-normality,
  finite-horizon transient gain, and Kreiss-constant diagnostics on $J_\ell^{(i)}$
  for the same models and prompts, and check whether they predict $\reff(\ell)$
  and its layer profile.
  \item Extend the certificate to nonlinear and sparse-autoencoder readouts,
  testing whether dark directions survive a stronger detector---the real safety
  question.
  \item Minimum-energy multi-layer control under an explicit collateral-damage
  constraint.
  \item A second-order certificate, extending validity past $\alpha^\star$ where
  most deployed steering operates.
  \item Benchmark certificate-guided design against circuit-discovery methods
  (ACDC, EAP) at equal model-run budget. \S\ref{sec:results} establishes the
  gain against random and probe-ordered selection, but not against those.
  \item A human-labelled toxicity behaviour functional, removing the
  logit-difference proxy from the safety claim.
\end{enumerate}

\section{Conclusion}

Reading a residual stream and writing to it are governed by different subspaces.
That separation is provable, constructible, and measurable in advance from one
backward pass---and on the models measured here it is close to total. Probe
accuracy, the field's default evidence that a direction ``is'' a concept,
carries essentially no information about whether that direction controls
anything.

\paragraph{Reproducibility.} Every number in this paper is generated from
\texttt{results/data/*.json} via \texttt{paper/make\_numbers.py}, and every
figure from the same files via \texttt{paper/make\_figures.py}. No measurement
is typed by hand into the text.

\appendix

\section{Proofs}
\label{app:proofs}

\begin{proof}[Proof of Theorem~\ref{thm:decomp}]
Taylor's theorem with Lagrange remainder on the segment gives
$\Delta_i = \alpha g^{(i)\top}w + R_i$ with
$|R_i|\le\frac12\alpha^2\|w\|^2\sup\|\nabla^2\Phi_\ell\|\le\frac{M_\ell}{2}\alpha^2\|w\|^2$
under A1. Splitting $w=P_Cw+(I-P_C)w$ is an identity. For the RMS bound,
$\frac1N\sum_i(g^{(i)\top}(I-P_C)w)^2 = w^\top(I-P_C)\Wg(I-P_C)w \le \lambda_{r+1}\|w\|^2$,
since $(I-P_C)\Wg(I-P_C)$ has spectral norm $\lambda_{r+1}$.
\end{proof}

\begin{proof}[Proof of Theorem~\ref{thm:decouple}]
Take $\Phi_\ell(x)=e_1^\top x$, so $g^{(i)}=\gbar=e_1$ for all $i$, and draw
$x\mid y\sim\mathcal{N}(y\Delta\mu\,e_2,\sigma^2 I_d)$. Then $\delta=\Delta\mu\,e_2$,
$\Sigma=\sigma^2I$, and $w_{\mathrm{probe}}\propto e_2$.
(b) $\Delta_i(\alpha,w_{\mathrm{probe}})=\Phi_\ell(x+\alpha e_2)-\Phi_\ell(x)=\alpha e_1^\top e_2=0$
exactly, for all $\alpha,i$.
(a) The score $e_2^\top x$ is $\mathcal{N}(y\Delta\mu,\sigma^2)$, so
$\mathrm{AUC}=\Phi_{\mathcal N}(\Delta\mu/(\sigma\sqrt2))$; choose
$\Delta\mu/\sigma \ge \sqrt2\,\Phi_{\mathcal N}^{-1}(1-\epsilon)$.
(c) Take $v=e_1$: its score is $\mathcal{N}(0,\sigma^2)$ independent of $y$, so
its AUC is exactly $\tfrac12$, while $|\gbar^\top v|=1=\|\gbar\|$ and
$\Delta_i(\alpha,v)=\alpha$.
\end{proof}

\begin{proof}[Proof of Proposition~\ref{prop:inert}]
$\mathbb{E}[ww^\top]=I/d$, so $\mathbb{E}[w^\top\Wg w]=\operatorname{tr}(\Wg)/d$.
The maximum of $w^\top\Wg w$ on the unit sphere is $\lambda_1$; divide.
\end{proof}

\begin{proof}[Proof of Proposition~\ref{prop:dark}]
Maximising $|\gbar^\top v|$ subject to $\|v\|=1$, $P_Ov=0$ is maximising
$|((I-P_O)\gbar)^\top v|$ over unit $v\in\mathcal{O}^\perp$, attained at the
normalised $(I-P_O)\gbar$. For the AUC claim, the score $v^\top x$ has
class-conditional means $v^\top\mu_y$, and $v^\top\delta$ vanishes when $v$ is
orthogonal to $(\Sigma+\gamma I)^{-1}\delta$ in the $(\Sigma+\gamma I)$-metric;
equal means with shared variance gives AUC $\tfrac12$.
\end{proof}

\begin{proof}[Proof sketch of Proposition~\ref{prop:leak}]
Write $\hat w = w_{\mathrm{pop}} + e$ with estimation error $e$ asymptotically
Gaussian and, for isotropic $\Sigma$, isotropic within the complement of
$w_{\mathrm{pop}}$ with $\mathbb{E}\|e\|^2 = Cd/N$. The leakage is the component
of $e$ along the single fixed unit vector $\gbar$, i.e.\ one coordinate of an
isotropic vector spread over $d-1$ dimensions:
$\mathbb{E}[(e^\top\gbar)^2] = \mathbb{E}\|e\|^2/(d-1) = Cd/(N(d-1)) \to C/N$.
The $d$ factors cancel.
\end{proof}

\section{Additional experimental detail}

\paragraph{Behaviour functional.} $\varphi$ is the mean logit over a set of mild
insult tokens minus the mean over a matched neutral set, at the final position.
Insult vocabulary is used strictly as a \emph{logit measurement target};
nothing is generated or emitted. Identity terms appear only to measure false
positives, which is their sole purpose.

\paragraph{Adjoints at every layer from one pass.} Forward pre-hooks capture the
input to each block and call \texttt{retain\_grad}; a single backward of
$\sum_i\varphi_i$ then yields $g_\ell^{(i)}$ at every layer simultaneously, so
the certificate costs one forward and one backward pass in total.

\paragraph{Curvature estimation.} $M_\ell$ is estimated by non-uniform second
differences of $\alpha\mapsto\overline{\Delta(\alpha,\gbar)}$ on a grid spanning
$[-2\alpha,2\alpha]$.

\begin{thebibliography}{99}
\small
\bibitem{eeao2025}
M.~Méloux, F.~Portet, S.~Maniu, and M.~Peyrard. \emph{Everything, Everywhere,
All at Once: Is Mechanistic Interpretability Identifiable?} ICLR, 2025.
arXiv:2502.20914.
\bibitem{uniq2026}
A.~Sahi, I.~Kutianawala, M.~Beeram, S.~Song, and A.~Shrivastava.
\emph{Characterizing Mechanistic Uniqueness and Identifiability Through
Circuit Analysis.} ICLR, 2026. OpenReview:cHpZ5DE8wj.
\bibitem{manycircuits2026}
A.~Bayat Makou, J.~Niu, S.~Dutta, and I.~Gurevych. \emph{Many Circuits, One
Mechanism: Input Variation and Evaluation Granularity in Circuit Discovery.}
2026. arXiv:2606.06267.
\bibitem{certified2026}
A.~Anani, T.~Lorenz, B.~Schiele, M.~Fritz, and J.~Fischer. \emph{Certified
Circuits: Stability Guarantees for Mechanistic Circuits.} ICML, 2026.
arXiv:2602.22968.
\bibitem{blackbox2025}
J.~Moon. \emph{From Black-Box to White-Box: Control-Theoretic Neural Network
Interpretability.} 2025. arXiv:2511.12852.
\bibitem{ssmctrl2025}
M.~Mabrok and Y.~Zafari. \emph{Controllability Analysis of State Space-based
Language Model.} 2025. arXiv:2511.17970.
\bibitem{langaxis2026}
A.~Srivastav. \emph{Steering the Language Axis: From Linear Decodability to
Causal Control.} 2026. arXiv:2608.12334.
\bibitem{actadd}
A.~M.~Turner, L.~Thiergart, G.~Leech, D.~Udell, U.~Mini, and M.~MacDiarmid.
\emph{Activation Addition: Steering Language Models Without Optimization.}
2024. arXiv:2308.10248.
\bibitem{caa}
N.~Rimsky, N.~Gabrieli, J.~Schulz, M.~Tong, E.~Hubinger, and A.~M.~Turner.
\emph{Steering Llama 2 via Contrastive Activation Addition.} ACL, 2024.
arXiv:2312.06681.
\bibitem{feedback2025}
D.~V.~Nguyen, H.~M.~Vu, N.~Y.~Pham, L.~Zhang, and T.~M.~Nguyen.
\emph{Activation Steering with a Feedback Controller.} ICLR, 2026.
arXiv:2510.04309.
\bibitem{locallin2026}
J.~Skifstad, X.~A.~Yang, and G.~Chou. \emph{Local Linearity of LLMs Enables
Activation Steering via Model-Based Linear Optimal Control.} 2026.
arXiv:2604.19018.
\bibitem{causaldetox2026}
Y.~Wang, Y.~Chen, A.~Goyal, and H.~Sundaram. \emph{CausalDetox: Causal Head
Selection and Intervention for Language Model Detoxification.} ACL, 2026.
arXiv:2604.14602.
\bibitem{kalman1963}
R.~E.~Kalman. \emph{Mathematical Description of Linear Dynamical Systems.}
SIAM Journal on Control, 1(2):152--192, 1963. DOI:10.1137/0301010.
\bibitem{backusgilbert1968}
G.~Backus and F.~Gilbert. \emph{The Resolving Power of Gross Earth Data.}
Geophysical Journal of the Royal Astronomical Society, 16(2):169--205, 1968.
DOI:10.1111/j.1365-246X.1968.tb00216.x.
\bibitem{elfving1952}
G.~Elfving. \emph{Optimum Allocation in Linear Regression Theory.} The
Annals of Mathematical Statistics, 23(2):255--262, 1952.
DOI:10.1214/aoms/1177729442.
\bibitem{trefethen1993}
L.~N.~Trefethen, A.~E.~Trefethen, S.~C.~Reddy, and T.~A.~Driscoll.
\emph{Hydrodynamic Stability Without Eigenvalues.} Science, 261(5121):578--584,
1993. DOI:10.1126/science.261.5121.578.
\bibitem{nonnormalresidual2026}
J.~Fernando and G.~Guitchounts. \emph{Dynamics of the Transformer Residual
Stream: Coupling Spectral Geometry to Network Topology.} 2026.
arXiv:2605.14258.
\bibitem{transientreserves2026}
H.~Li. \emph{Transient Reserves, Sink Dampers, and the Failure of Eigenvalue
Reasoning in the Attention Propagator.} 2026. arXiv:2607.09279.
\end{thebibliography}

\bigskip
\noindent\fbox{\begin{minipage}{0.96\linewidth}
\small
\textbf{Bibliography status.} All entries above were verified against
primary sources (arXiv, ACL Anthology, ICML/ICLR listings, OpenReview,
Project Euclid, Science) on 2026-08-26; author lists, venues, and
identifiers are as confirmed at that date. Four references beyond the
original twelve---Kalman (1963), Backus and Gilbert (1968), Elfving (1952),
and Trefethen et al. (1993)---ground the structural connections drawn in
\S\ref{sec:two}, \S\ref{sec:results}, and Related Work. Two further 2026
preprints report direct empirical measurements of the Transformer-Jacobian
non-normality discussed in Remark~\ref{rem:nonnormal}; as very recent,
single-venue preprints they are flagged as concurrent evidence, not
independently replicated results.
\end{minipage}}

\end{document}
